\documentclass[11pt,a4paper]{article}
\usepackage[UTF8]{ctex}
\usepackage[left=2.2cm,right=2.2cm,top=2.4cm,bottom=2.4cm]{geometry}
\usepackage{amsmath,amssymb}
\usepackage{booktabs}
\usepackage{longtable}
\usepackage{array}
\usepackage{enumitem}
\usepackage{hyperref}
\usepackage{cite}

% bstctlcite support under article class
\makeatletter
\providecommand{\bstctlcite}[1]{%
  \@bsphack
  \@for\@citeb:=#1\do{%
    \edef\@citeb{\expandafter\@firstofone\@citeb}%
    \if@filesw\immediate\write\@auxout{\string\citation{\@citeb}}\fi
  }%
  \@esphack
}
\makeatother

\title{FR3 力控 GRU 仿真闭环迁移实施文档\\
\large Gazebo + ROS2 + C++ with Hybrid 兼容迁移}
\author{Mingdao Lin}
\date{2026-03-04}

\begin{document}
\maketitle
\bstctlcite{BSTcontrol}

\section{目标与范围}
本文定义 Franka Research 3 平台上的力控迁移计划，目标是完成完整闭环，即数据采集 -- 训练 -- 仿真测试 -- 统计验证 -- 失败复盘。本文只覆盖门控循环单元，Gated Recurrent Unit，GRU，暂不引入 physics-guided 结构。系统主线采用直接关节力矩，direct joint torque，控制接口，反馈主稳定器为比例积分微分，Proportional-Integral-Derivative，PID，前馈项由 GRU 产生。

当前资产包括 TCST 训练脚本与结果目录，见 \texttt{Codes\_Setup/Trainer\_GRU.py} 与 \texttt{Training\_results/GRU}。现有脚本是单输入单输出，Single-Input Single-Output，SISO，输入 \texttt{yd\_filtered}，输出 \texttt{u}。迁移策略采用 Hybrid 兼容迁移，即旧脚本保持可运行，新数据通过适配层自动投影到旧格式，同时保留新格式字段用于 MIMO 扩展。

\section{数学定义}
\subsection{变量与控制律}
接触坐标系期望力，desired wrench，记为 $w_d[k]$，实测力，measured wrench，记为 $w_m[k]$，误差定义为
\begin{equation}
  e_w[k] = w_d[k] - w_m[k].
\end{equation}

关节角与关节角速度分别记为 $q[k] \in \mathbb{R}^7$，$\dot q[k] \in \mathbb{R}^7$。任务雅可比，task Jacobian，记为 $J_t(q[k])$。PID 力反馈力矩定义为
\begin{equation}
  \tau_{pid}[k] = J_t(q[k])^T \left(K_p e_w[k] + K_i I_w[k] + K_d \dot e_w[k]\right).
\end{equation}

GRU 前馈定义为
\begin{align}
  h[k] &= \mathrm{GRU}\left(h[k-1], x[k]\right), \\
  w_{ff}[k] &= W_o h[k] + b_o, \\
  \tau_{ff}[k] &= J_t(q[k])^T w_{ff}[k].
\end{align}

最终命令定义为
\begin{equation}
  \tau_{cmd}[k] = \mathrm{sat}\left(\tau_{pid}[k] + \tau_{ff}[k]\right).
\end{equation}

\subsection{Phase 迁移}
\begin{itemize}[leftmargin=1.2em]
  \item Phase-1: $w_{ff}[k]=[0,0,\Delta F_z[k]]^T$，先跑通 Fz 单通道。
  \item Phase-2: $w_{ff}[k]=[\Delta F_x[k],\Delta F_y[k],\Delta F_z[k]]^T$，扩展三轴力。
  \item Phase-3: $w_{ff}[k]\in\mathbb{R}^6$，扩展到全 wrench。
\end{itemize}

\subsection{损失函数与安全边界}
训练损失定义为
\begin{equation}
\begin{aligned}
L = &\lambda_1 L_{track} + \lambda_2 L_{supervised} + \lambda_3 L_{smooth} \\
    &+ \lambda_4 L_{limit} + \lambda_5 L_{rate},
\end{aligned}
\end{equation}
其中
\begin{align}
L_{track} &= \mathrm{mean}\left(\|w_d[k+1]-w_m[k+1]\|^2\right), \\
L_{supervised} &= \mathrm{mean}\left(\|\tau_{ff}^{pred}[k]-\tau_{ff}^{target}[k]\|^2\right), \\
L_{smooth} &= \mathrm{mean}\left(\|\tau_{ff}[k]-\tau_{ff}[k-1]\|^2\right), \\
L_{limit} &= \mathrm{mean}\left(\|\max(0,|\tau_{cmd}|-\tau_{max})\|^2\right), \\
L_{rate} &= \mathrm{mean}\left(\|\max(0,|\Delta\tau_{cmd}|-\Delta\tau_{max})\|^2\right).
\end{align}

在线边界为
\begin{align}
|\tau_{cmd,i}| &\le \tau_{max,i}, \\
|\Delta\tau_{cmd,i}| &\le \Delta\tau_{max,i}, \\
|\tau_{ff,i}| &\le \alpha |\tau_{pid,i}| + \beta_i.
\end{align}

当出现传感器异常 such as topic timeout、接触丢失 such as \texttt{contact\_flag=0}、误差爆发 such as $|e_w|$ 连续超阈值时，强制 $\tau_{ff}=0$ 并回退纯 PID。

\section{工程实现}
\subsection{ROS2 与 C++}
已实现新包 \texttt{fr3\_gru\_force}，核心节点如下。
\begin{itemize}[leftmargin=1.2em]
  \item \texttt{gru\_force\_controller\_node}: 订阅 \texttt{/fr3/state}、\texttt{/fr3/wrench\_t}、\texttt{/fr3/contact\_flag}，发布 \texttt{/fr3/torque\_cmd\_gru}，并发布调试 topic \texttt{/gru/debug/w\_ff}、\texttt{/gru/debug/tau\_ff}、\texttt{/gru/debug/tau\_pid}、\texttt{/gru/debug/safety\_flag}。
  \item \texttt{trial\_logger\_node}: 记录扩展 CSV schema，覆盖 \texttt{tau\_pid\_0..6}、\texttt{tau\_ff\_target\_0..6}、\texttt{tau\_cmd\_0..6}、\texttt{w\_ref\_fx..fz}、\texttt{w\_meas\_fx..fz}、\texttt{sync\_err\_ms}、\texttt{delay\_est\_ms}。
\end{itemize}

新增 launch 文件 \texttt{fr3\_gazebo\_gru.launch.py}，参数文件 \texttt{gru\_params.yaml}。

\subsection{Python pipeline}
已实现脚本如下。
\begin{itemize}[leftmargin=1.2em]
  \item \texttt{tcst\_hybrid\_adapter.py}: FR3 统一 schema 与 TCST 文件夹双向转换。
  \item \texttt{train\_gru\_force.py}: Phase-1 训练，输出 \texttt{model.pt}、\texttt{norm\_stats.npz}、\texttt{model\_meta.yaml}、\texttt{train\_report.yaml}。
  \item \texttt{export\_gru\_onnx.py}: 导出 \texttt{model.onnx}。
  \item \texttt{evaluate\_gru\_force.py}: 基线对比与统计检验，输出 \texttt{gru\_phase1\_eval.yaml}。
  \item \texttt{run\_gru\_phase1\_loop.py}: 一键执行 train -- export -- evaluate。
\end{itemize}

\section{默认参数与实验矩阵}
\subsection{采样与激励}
默认采样率为 \texttt{physics 1000 Hz} 与 \texttt{control 200 Hz}。激励覆盖 chirp、PRBS、multisine 与切向轨迹簇，范围与离散点在 \texttt{sim/contact\_envs.yaml} 中定义。环境刚度包含 ID 集合 \texttt{500, 5000, 50000 N/m} 与 OOD 集合 \texttt{1500, 15000 N/m}。

\subsection{基线与指标}
比较三条基线。
\begin{itemize}[leftmargin=1.2em]
  \item PID only
  \item linear FF + PID
  \item GRU FF + PID
\end{itemize}

指标为 IAE、RMSE、峰值误差、响应时间、控制能量、约束违规率 CVR。统计过程采用 Friedman 总体检验与 Wilcoxon signed-rank 事后检验，Holm 修正多重比较。

\section{Conda 与部署建议}
训练环境推荐 \texttt{configs/conda/fr3-train.yml}，使用 Python 3.10 与 PyTorch 2.5.1。ROS2 运行与训练环境隔离，建议 ROS2 使用系统 Python 3.10，工具侧可选 \texttt{configs/conda/fr3-ros-tools.yml}。该隔离策略避免 \texttt{PYTHONPATH} 与 \texttt{LD\_LIBRARY\_PATH} 交叉污染。

\section{里程碑与验收}
\begin{longtable}{p{0.16\linewidth}p{0.36\linewidth}p{0.40\linewidth}}
\toprule
阶段 & 交付件 & 验收条件 \\
\midrule
WP0 & conda 与 ROS2 预装验收 & 训练脚本导入通过，ROS2 包编译通过 \\
WP1 & 直接力矩链路 + watchdog & topic 连通，异常触发一周期内回退 \\
WP2 & 分层采集日志 & CSV 列完整，时间同步误差可追踪 \\
WP3 & Hybrid 适配层 & FR3 转 TCST 与 TCST 转 FR3 双向可用 \\
WP4 & Fz GRU 训练 & 产出 \texttt{model.pt} 与 \texttt{train\_report.yaml} \\
WP5 & 仿真闭环评估 & 三基线评估与统计检验报告可复现 \\
WP6 & FxFyFz 扩展草案 & 保持同一评估矩阵与安全边界 \\
WP7 & 文档与引用核验 & 单栏中文 LaTeX 编译成功，DOI 全部可达 \\
\bottomrule
\end{longtable}

\section{文献与工程基线}
本文引用近 15 年文献，覆盖 TCST、TAC、IJRR、TRO、RA-L 与工程实现仓库。关键方法来源包括 inversion-based feedforward、hybrid force-impedance、control-oriented meta-learning 与 variable impedance learning。对应 DOI 和 URL 见参考文献条目\cite{lin2026pgg,xu2025hybridff,li2014annff,spiegel2024stableinversion,haddadin2024ufic,richards2023coml,buchli2011vic,li2025tro,iskandar2023ral,franka_ros2,franka_ign,franka_gz_sim}。

\bibliographystyle{IEEEtran}
\bibliography{references}

\end{document}
